\documentclass[12pt,-letter paper]{article}
\usepackage{tfrupee}
\usepackage{textcomp}
\usepackage{siunitx}
\providecommand{\sbrak}[1]{\ensuremath{{}\left[#1\right]}}
\usepackage{setspace}
\usepackage{gensymb}
\usepackage{xcolor}
\usepackage{caption}
%\usepackage{subcaption}
\doublespacing
\singlespacing
\usepackage[none]{hyphenat}
\usepackage{amssymb}
\usepackage{relsize}
\usepackage[cmex10]{amsmath}
\usepackage{mathtools}
\usepackage{amsmath}
\usepackage{commath}
\usepackage{amsthm}
\interdisplaylinepenalty=2500
%\savesymbol{iint}
\usepackage{txfonts}
%\restoresymbol{TXF}{iint}
\usepackage{wasysym}
\usepackage{amsthm}
\usepackage{mathrsfs}
\usepackage{txfonts}
\let\vec\mathbf{}
\usepackage{stfloats}
\usepackage{float}
\usepackage{cite}
\usepackage{cases}
\usepackage{subfig}
%\usepackage{xtab}
\usepackage{longtable}
\usepackage{multirow}
%\usepackage{algorithm}
\usepackage{amssymb}
%\usepackage{algpseudocode}
\usepackage{enumitem}
\usepackage{mathtools}
%\usepackage{eenrc}
%\usepackage[framemethod=tikz]{mdframed}
\usepackage{listings}
%\usepackage{listings}
\usepackage[latin1]{inputenc}
%%\usepackage{color}{   
%%\usepackage{lscape}
\usepackage{textcomp}
\usepackage{titling}
\usepackage{hyperref}
%\usepackage{fulbigskip}   
\usepackage{tikz}
\usepackage{graphicx}
\lstset{
  frame=single,
  breaklines=true
}
\let\vec\mathbf{}
\usepackage{enumitem}
\usepackage{graphicx}
\usepackage{siunitx}
\let\vec\mathbf{}
\usepackage{enumitem}
\usepackage{graphicx}
\usepackage{enumitem}
\usepackage{tfrupee}
\usepackage{amsmath}
\usepackage{amssymb}
\usepackage{mwe} % for blindtext and example-image-a in example
\usepackage{wrapfig}
\graphicspath{{figs/}}
\providecommand{\mydet}[1]{\ensuremath{\begin{vmatrix}#1\end{vmatrix}}}
\providecommand{\myvec}[1]{\ensuremath{\begin{bmatrix}#1\end{bmatrix}}}
\providecommand{\cbrak}[1]{\ensuremath{\left\{#1\right\}}}
\providecommand{\brak}[1]{\ensuremath{\left(#1\right)}}
\begin{document}                                
\begin{center}                                                               
\textbf{MATRIX} 
\end{center} 
\begin{enumerate}
	\item If the matrix \(\myvec{6 & -x^2 \\ 2x-15 & 10}\)is symmetrix,find the value of $x$.
	\item Using properties of determinants,prove that: \\
		$\begin{vmatrix} a & b & b+c \\ c & b & c+a \\ b& c & a+b \end{vmatrix}$=\brak{a+b+c}$\brak{a-c}^2$
		\item Given that: $A$=\(\myvec{1 & -1 & 0 \\ 2 & 3 & 4 \\ 0 & 1 & 2}\)and $B$=\(\myvec{2 & 2 & -4 \\ -4 & 2 & -4 \\ 2 & -1 & 5}\). find $AB$. \\
			Using this result,solve the following system of equation: \\
			$ x-y=3 $,$2x+3y+4z=17$ and $y+2z=7$
\end{enumerate}
\begin{center}                                                               
\textbf{Geometry}
\end{center}
    \begin{enumerate}     
	    \item If $ y-2x-k $= $0 $ touches the conic $ 3x^2-5y^2$ = $15$, find the value of $k$.  
             \item Show that the lines $\frac{x-4}{1}$=$\frac{y+3}{-4}$=$\frac{z+1}{7}$ and $\frac{x-1}{2}$=$\frac{y+1}{-3}$=$\frac{z+10}{8}$ intersect. Find the coordinates of their point of intersection.
     \item Find the equation of the plane passing through the point  \brak{1,-2, 1} and perpendicular to the line joining the points A \brak{3, 2, 1} and B\brak{1, 4, 2}.                                                   
\end{enumerate} 
\begin{center}
	\textbf{Trignometry}
\end{center}
\begin{enumerate}
	\item Prove that$\frac{1}{2}$ $\cos^{-1}$ \brak{\frac{1-x}{1+x}}=$\tan^{-1}$ $\sqrt{x}$ 	
	\item Solve the equation for $x$: \\
		$\sin^{-1}$$x$+$\sin^{-1}$$\brak{1-x}$=$\cos^{-1}$$x$,$x\neq 0$
\end{enumerate}
\begin{center}
	\textbf{limits and continuity}
\end{center}
\begin{enumerate}
	\item Using L'Hospital's Rule,evaluate: \\
		$\lim_{x\to\frac{\pi}{2}}\brak{x\tan{x}-\frac{\pi}{4}.\sec{x}}$
\end{enumerate}
\begin{center}
	\textbf{Integrations}
\end{center}
\begin{enumerate}
	\item Evaluate: $\int \frac{1}{x^2} \sin^2 \brak{\frac{1}{x}} dx $
	\item Evaluate: \\
		$\int_{0}^\frac{\pi}{4}$ log\brak{1+\tan{\theta}} $d\theta$
	\item Evaluate: \\
		$\int \frac{\sin2x} {\brak{1+\sin{x}}\brak{2+\sin{x}}}dx$
\end{enumerate}
\begin{center}
	\textbf{statistics}
\end{center}
\begin{enumerate}
	\item By using the data $\overline{x}$=$25$,$\overline{y}$=$30$,$b_{yx}$=$1.6$ and $b_{xy}$=$0.4$,find:
		\begin{enumerate}
	\item The regression equation $y$ on $x$. 
  \item What is the most likely value of $y$ when $x$ = $60?$ 
  \item What is the coefficient of correlation between $x$ and $y?$
\end{enumerate}
	\item Calculate the Spearman's rank correlation coefficient for the following data and interpret the result: \\
		\begin{tabular}{|c|c|c|c|c|c|c|c|c|c|c|}
	\hline
			X & 35 & 54 & 80 & 95 & 73 & 73 & 35 & 91 & 83 & 81\\
  \hline				       	       Y & 40 & 60 & 75 & 90 & 70 & 75 & 38 & 95 & 75 & 70\\
					       \hline
					              \end{tabular}
	\item Find the line of best fit for the following data, treating $x$ as dependent variable \\
		Regression equation $x$ on $y$: \\
\begin{tabular}{|c|c|c|c|c|c|c|c|c|}
			\hline
	X & 14 & 12 & 13 & 14  & 16 & 10 & 13 & 12\\
	\hline
Y & 14 & 23 & 17 & 24 & 18 & 25 & 23 & 24\\		\hline
					              \end{tabular} \\
						      Hence,estimate the value of $x$ when$y=16$.
					      \item The price relatives and weights of a set of commodities are given below: \\
						       \begin{tabular}{|c|c|c|c|c|}
							       \hline
	Commodity& A & B & C & D\\
								       \hline		
							Price relatives & 125 & 120 & 127 & 119\\
									       \hline
									Weights & x & 2x &y & y+3\\
									\hline
	        \end{tabular} \\
If the sum of the weights is $40$ and weighted average of price relatives index number is $122$, find the numerical values of $x$ and $y$.
\item Construct $3$ yearly moving averages from the following data and show on a graph against the original data: \\
	 \begin{tabular}{|c|c|c|c|c|c|c|c|c|c|c|}
		 \hline
		 	       Year& 2000 & 2001 & 2002 & 2003 & 2004 & 2005 & 2006 & 2007 & 2008 & 2009\\
			       \hline
 Annualsalesinlakhs & 18 & 22 & 20 & 26 & 30 & 22 & 24 & 28 & 32 & 35\\
				       \hline
				              \end{tabular}

\end{enumerate}
\begin{center}
	\textbf{probability}
\end{center}
\begin{enumerate}
		 \item A problem is given to three students whose chances of solving it are $\frac{1}{4}$,$\frac{1}{5}$ and $\frac{1}{3}$ respectively. Find the probability that the problem is solved.
		 \item In a class of $60$ students,$30$ opted for Mathematics, $32$ opted for Biology and $24$ opted for both Mathematics and Biology. If one of these students is selected at random, find the probability that:  \begin{enumerate}
				 \item The student opted for Mathematics or Biology. 
				 \item The student has opted neither Mathematics nor Biology. 
		 \item The student has opted Mathematics but not Biology. 
		 \end{enumerate}
	 \item Bag $A$ contains $1$ white,$ 2$ blue and $3$ red balls. Bag $B$ contains $3$ white, $3$ blue and $2$ red balls. Bag $C$ contains $2$ white, $3$ blue and $4$ red balls. One bag is selected at random and then two balls are drawn from the selected bag. Find the probability that the balls drawn are white and red.
	 \item A fair die is rolled. If face $1$ turns up, a ball is drawn from Bag $A$. If face $2$ or $3$ turns up, a ball is drawn from Bag $B$. If face $4$ or $5$ or $6$ turns up, a ball is drawn from Bag $C$. Bag $A$ contains $3$ red and $2$ white balls, Bag $B$ contains $3$ red and $4$ white balls and Bag $C$ contains $4$ red and $5$ white balls. The die is rolled, a Bag is picked up and a ball is drawn. If the drawn ball is red, what is the probability that it is drawn from Bag $B?$ 
	 \item An urn contains $25$ balls of which $10$ balls are red and the remaining green. A ball is drawn at random from the urn, the colour is noted and the ball is replaced. If $6$ balls are drawn in this way, find the probability that: 
		 \begin{enumerate}
		\item All the balls are red. 
		\item Not more than $2$ balls are green. 
		\item Number of red balls and green balls are equal.
		 \end{enumerate}

\end{enumerate}
\begin{center}
	\textbf{complexnumbers}
\end{center}
\begin{enumerate}
	\item If $a+ib$=$\frac{x+iy}{x-iy}$ Prove that $a^2+b^2$=$1$ and $\frac{b}{a}$=$\frac{2xy}{x^2-y^2}$
	\item Prove that locus of $z$ is circle and find its centre and radius if $\frac{z-i}{z-1}$ is purely imaginary. 
\end{enumerate}
\begin{center}
	\textbf{Diffrentiations}
\end{center}
\begin{enumerate}
	\item Solve: $\frac{dy}{dx}$=$1-xy+y-x$
	\item If $y$=$\cos$\brak{\sin{x}},show that: \\
		$\frac{d^2y}{dx^2}$+$\tan{x}$$\frac{dy}{dx}$+$y$$\cos^2${x}=$0$
	\item Solve:  $\brak{x^2-yx^2}$$dy$+$\brak{y^2+xy^2}$$dx$=$0$
\end{enumerate}
\begin{center}
	\textbf{BooleanAlgebra}
\end{center}
\begin{enumerate}
	\item  If $A$, $B$ and $C$ are the elements of Boolean algebra, simplify the expression $\brak{A'+B'}$$\brak{A + C'}$ +$ B'$$\brak{B + C}$. Draw the simplified circuit.
\end{enumerate}
\begin{center}
	\textbf{Derivatives}
\end{center}
	\begin{enumerate}
		\item Verify Langrange's mean value theorem for the function: \\ 
$f$\brak{x} = $x$ $\brak{1 - log x}$ and find the value of $'c'$ in the interval \sbrak{1,2} 
\item Show that the surface area of a closed cuboid with square base and given volume is minimum when it is a cube.
	\end{enumerate}
	\begin{center}
		\textbf{Hyperbola}
	\end{center}
	\begin{enumerate}
		\item Find the coordinates of the centre, foci and equation of directrix of the hyperbola $x^2-3y^2-4x$=$8$.
	\end{enumerate}
	\begin{center}
		\textbf{Parabola}
	\end{center}
	\begin{enumerate}
		\item Draw a rough sketch of the curve $y^2$= $4x$ and find the area of the region enclosed by the curve and the line $y$=$x$.
	\end{enumerate}
\begin{center}
	\textbf{Vectors}
\end{center}
\begin{enumerate}
  \item If $\overrightarrow{a}$, $\overrightarrow{b}$, $\overrightarrow{c}$ are three mutually perpendicular vectors of equal magnitude, prove that $\brak{\overrightarrow{a} + \overrightarrow{b} + \overrightarrow{c}}$ is equally inclined with vectors $\overrightarrow{a}$, $\overrightarrow{b}$, and $\overrightarrow{c}$.
  \item Find the value of $\lambda$ for which the four points with position vectors $6\hat{i}-7\hat{j}$,$16\hat{i}-19\hat{j}-4\hat{k}$,$\lambda\hat{j}-6\hat{k}$ and $2\hat{i}-5\hat{j}+10\hat{k}$ are coplanar.
\end{enumerate}
\begin{center}
	\textbf{profitandloss}
\end{center}
\begin{enumerate}
	\item A machine costs \rupee $60,000$ and its effective life is estimated to be $25$ years. A sinking fund is to be created for replacing the machine at the end of its life time when its scrap value is estimated as \rupee $5,000$. The price of the new machine is estimated to be $100$\% more than the price of the present one. Find the amount that should be set aside at the end of each year, out of the profits, for the sinking fund if it accumulates at an interest of $6$\% per annum compounded annually. 
	\item A farmer has a supply of chemical fertilizer of type $A$ which contains $10$\% nitrogen and $6$\% phosphoric acid and of type $B$ which contains $5$\% nitrogen and $10$\% phosphoric acid. After soil test, it is found that at least $7$ kg of nitrogen and same quantity of phosphoric acid is required for a good crop. The fertilizer of type $A$ costs \rupee $ 5.00$ per kg and the type $B$ costs \rupee $8.00$ per kg. Using Linear programming, find how many kilograms of each type of the fertilizer should be bought to meet the requirement and for the cost to be minimum. Find the feasible region in the graph.
	\item The demand for a certain product is represented by the equation$p$=$500+25x-\frac{x^2}{3}$in rupees where $x$ is the number of units and $p$ is the price per unit. Find: 
		\begin{enumerate}
\item Marginal revenue function. 
\item  The marginal revenue when $10$ units are sold.
		\end{enumerate}
	\item A bill of \rupee $60,000$ payable $10$ months after date was discounted for \rupee $ 57,300$ on $30^{th}$ June,$2007$. If the rate of interest was $11$$\frac{1}{4}$\% per annum, on what date was the bill drawn?

\end{enumerate}
\end{document}

